% Table 4: 0-shot downstream (lm-eval-harness). lambada: acc (ppl).
\begin{table*}[t]\centering\footnotesize
\begin{tabular}{lccccccc}
\toprule
Model & lambada (ppl) & sciq & arc\_easy & piqa & hellaswag & winogrande & mmlu \\
\midrule
chance & $\sim$0 & 0.25 & 0.25 & 0.50 & 0.25 & 0.50 & 0.25 \\
\midrule
SAN 24M @31B & 0.091 (821) & 0.725 & 0.397 & 0.545 & 0.279 & 0.504 & 0.236 \\
FFN iso-param @31B & 0.136 (492) & 0.702 & 0.416 & 0.563 & 0.283 & 0.516 & 0.233 \\
SAN 24M @105B (3s) & 0.111 (682) & 0.742 & 0.405 & 0.559 & 0.284 & 0.517 & 0.232 \\
FFN iso-param @105B (3s) & 0.134 (518) & 0.661 & 0.415 & 0.558 & 0.283 & 0.511 & 0.245 \\
FFN iso-FLOP 43M @105B & 0.178 (451) & 0.634 & 0.419 & 0.562 & 0.293 & 0.534 & 0.230 \\
FFN iso-depth 87M @105B & 0.111 (964) & 0.630 & 0.426 & 0.552 & 0.288 & 0.525 & 0.233 \\
\midrule
SAN fineweb-trained @31B & 0.203 (189) & 0.705 & 0.447 & 0.579 & 0.294 & 0.507 & 0.231 \\
FFN fineweb-trained @31B & 0.181 (218) & 0.707 & 0.478 & 0.588 & 0.287 & 0.507 & 0.232 \\
\bottomrule
\end{tabular}
\caption{The storage/routing split at task level. Lambada (out-of-distribution
recall for SYNTH-trained models) favors the FFN at every budget, but is
non-monotone in FFN capacity: iso-depth (87M), the best in-distribution model,
is the worst at it (0.111, ppl 964). Sciq (support passage in context) favors
the SAN with a margin that grows with training in the pre-registered direction
(seed ranges non-overlapping at 105B) and falls monotonically with FFN
capacity. Trained on knowledge-dense fineweb-edu, the lambada preference
reverses: a task's storage/routing identity is relative to the training
distribution. Remaining tasks are at chance at these scales.}
\label{tab:downstream}
\end{table*}
